\documentclass[10pt]{article}
\providecommand{\FIGDIR}{fig}
\newcommand{\DIGESTDATE}{16 August 2026}
% ===================== PM Research Watch - shared preamble =====================
\usepackage[T1]{fontenc}
\usepackage[utf8]{inputenc}
\usepackage{amsmath,amssymb}
\usepackage{tgpagella}
\usepackage{tgheros}
\renewcommand{\sfdefault}{qhv}
\renewcommand{\ttdefault}{lmtt}
\usepackage{microtype}
\usepackage[english]{babel}

\usepackage[a4paper,top=17mm,bottom=17mm,left=17mm,right=17mm,
            headsep=6mm,footskip=10mm]{geometry}
\usepackage{graphicx}
\usepackage{xcolor}
\usepackage{colortbl}
\usepackage{tikz}
\usetikzlibrary{calc}
\usepackage[most]{tcolorbox}
\usepackage{titlesec}
\usepackage{fancyhdr}
\usepackage{booktabs}
\usepackage{array}
\usepackage{longtable}
\usepackage{enumitem}
\usepackage{multicol}
\usepackage{ragged2e}
\usepackage[hidelinks]{hyperref}

% ------------------------------------------------------------------ palette
\definecolor{Ink}    {HTML}{14262E}
\definecolor{Deep}   {HTML}{0F4C5C}
\definecolor{Teal}   {HTML}{1F7A8C}
\definecolor{Sky}    {HTML}{4FA3C4}
\definecolor{Sage}   {HTML}{6FA287}
\definecolor{Amber}  {HTML}{E3A008}
\definecolor{Clay}   {HTML}{B5643C}
\definecolor{Coral}  {HTML}{D1495B}
\definecolor{Violet} {HTML}{6A5B9E}
\definecolor{Slate}  {HTML}{4A6572}
\definecolor{Paper}  {HTML}{FBF9F5}
\definecolor{Mist}   {HTML}{EEEAE2}
\definecolor{Rule}   {HTML}{D9D3C9}

\pagecolor{Paper}
\color{Ink}

\hypersetup{colorlinks=true, linkcolor=Deep, urlcolor=Teal, citecolor=Teal,
            pdftitle={PM Research Watch}, pdfauthor={Automated daily digest}}

% ------------------------------------------------------------------ headings
\titleformat{\section}
  {\sffamily\bfseries\large\color{Deep}}{}{0pt}
  {}[{\vspace{1.2mm}\color{Amber}\titlerule[1.1pt]}]
\titlespacing*{\section}{0pt}{5.5mm}{2.8mm}

\titleformat{\subsection}
  {\sffamily\bfseries\normalsize\color{Teal}}{}{0pt}{}
\titlespacing*{\subsection}{0pt}{3.6mm}{1.4mm}

\setlist[itemize]{leftmargin=4.2mm, itemsep=1.1mm, topsep=1.3mm, parsep=0pt,
                  label=\textcolor{Teal}{\footnotesize$\blacktriangleright$}}

% ------------------------------------------------------------------ header/footer
\pagestyle{fancy}
\fancyhf{}
\renewcommand{\headrulewidth}{0pt}
\renewcommand{\footrulewidth}{0pt}
\fancyhead[L]{\sffamily\scriptsize\color{Slate}PM RESEARCH WATCH\;\textcolor{Rule}{\textbar}\;\DIGESTDATE}
\fancyhead[R]{\sffamily\scriptsize\color{Slate}Particulate matter monitoring \& health effects}
\fancyfoot[C]{\sffamily\scriptsize\color{Slate}\thepage}

% ------------------------------------------------------------------ components
% metric tile
\newcommand{\metric}[3]{%
  \begin{tcolorbox}[enhanced, colback=white, colframe=Rule, boxrule=0.5pt,
    arc=1.6mm, left=2.2mm, right=2.2mm, top=1.8mm, bottom=1.8mm,
    borderline west={1.7mm}{0pt}{#3}, nobeforeafter, valign=center,
    height=17mm, width=\linewidth]
    \RaggedRight
    {\sffamily\bfseries\fontsize{17}{17}\selectfont\color{#3}#1}\\[0.5mm]
    {\sffamily\fontsize{6.5}{7.4}\selectfont\color{Slate}#2}
  \end{tcolorbox}}

% pull-quote / key insight box
\newtcolorbox{keybox}[1][]{enhanced, colback=Mist, colframe=Mist,
  arc=1.6mm, left=3.2mm, right=3.2mm, top=2.4mm, bottom=2.4mm,
  borderline west={1.7mm}{0pt}{Amber}, #1}

% subtopic band
\newcommand{\band}[2]{%
  \vspace{3.2mm}\noindent
  \begin{tcolorbox}[enhanced, colback=#2, colframe=#2, arc=1.2mm,
    left=2.6mm, right=2.6mm, top=1.1mm, bottom=1.1mm, boxrule=0pt,
    width=\linewidth, nobeforeafter]
    {\sffamily\bfseries\small\color{white}#1}
  \end{tcolorbox}\vspace{1.6mm}}

% one paper entry: {short cite}{title}{journal}{doi}{body}
\newcommand{\paperentry}[5]{%
  \par\vspace{1.9mm}\noindent
  {\sffamily\bfseries\footnotesize\color{Ink}#1}\;%
  {\sffamily\fontsize{7}{8}\selectfont\color{Slate}#3}\\[0.4mm]
  {\footnotesize\itshape\color{Deep}#2}\\[0.6mm]
  {\footnotesize #4}\\[0.3mm]
  {\fontsize{6.6}{7.6}\selectfont\color{Slate}\texttt{doi:}\,\href{https://doi.org/#5}{\texttt{\detokenize{#5}}}}
  \par}

% caption style
\newcommand{\figcap}[1]{\par\vspace{1.0mm}
  {\fontsize{7.2}{8.6}\selectfont\color{Slate}\textsf{\textbf{Fig.}}\; #1\par}\vspace{1.5mm}}

\providecommand{\permil}{\textperthousand}

\setlength{\parindent}{0pt}
\setlength{\parskip}{0pt}
\renewcommand{\arraystretch}{1.18}


\begin{document}

% ==================================================================== MASTHEAD
\thispagestyle{empty}
\begin{tikzpicture}[remember picture, overlay]
  \fill[Deep] (current page.north west) rectangle
      ($(current page.north east)+(0,-67mm)$);
  \fill[Amber] ($(current page.north west)+(0,-67mm)$) rectangle
      ($(current page.north east)+(0,-68.4mm)$);
  \foreach \x/\y/\r in {22/12/0.9, 41/26/1.5, 63/9/0.7, 88/31/1.1, 112/17/1.9,
                        138/29/0.8, 159/13/1.3, 178/54/1.0, 196/21/1.6,
                        30/44/0.6, 74/57/1.2, 125/48/0.9, 168/6/0.8, 191/59/0.7,
                        52/18/1.0, 100/60/0.6, 148/20/1.1, 66/49/1.4, 205/45/0.8}
    \fill[Sky, opacity=0.30] ($(current page.north west)+(\x mm,-\y mm)$) circle (\r mm);
\end{tikzpicture}

\vspace*{4mm}
{\sffamily\fontsize{7.6}{9}\selectfont\color{Sky}%
 AUTOMATED DAILY LITERATURE DIGEST\;\textbullet\;PhD RESEARCH WATCH\par}
\vspace{2.2mm}
{\sffamily\bfseries\fontsize{27}{29}\selectfont\color{white}%
 Particulate Matter\par}
{\sffamily\fontsize{27}{29}\selectfont\color{Sky}%
 Monitoring \& Health Effects\par}
\vspace{5.0mm}
{\sffamily\fontsize{8.4}{10}\selectfont\color{white}%
 Issue 2026-08-16\;\;\textcolor{Amber}{\textbar}\;\;Entry window 16 August 2026\;\;%
 \textcolor{Amber}{\textbar}\;\;10 records carried, 15 rejected\par}
\vspace{18mm}

% -------------------------------------------------------------- metric strip
\noindent
\begin{minipage}[t]{0.190\linewidth}\metric{10}{RECORDS CARRIED\\(15 REJECTED)}{Deep}\end{minipage}\hfill
\begin{minipage}[t]{0.190\linewidth}\metric{6}{SUBTOPIC CLUSTERS\\REPRESENTED}{Teal}\end{minipage}\hfill
\begin{minipage}[t]{0.190\linewidth}\metric{11}{EFFECT ESTIMATES\\WITH INTERVALS}{Sky}\end{minipage}\hfill
\begin{minipage}[t]{0.190\linewidth}\metric{5}{INSTRUMENTATION /\\EXPOSURE RECORDS}{Amber}\end{minipage}\hfill
\begin{minipage}[t]{0.190\linewidth}\metric{2}{PREPRINTS\\(LABELLED)}{Clay}\end{minipage}

\vspace{6mm}

% -------------------------------------------------------------- signal of day
\section*{\sffamily\bfseries\color{Deep}Signal of the day}
\vspace{-1.5mm}{\color{Amber}\rule{\linewidth}{1.1pt}}\vspace{2.5mm}

\begin{keybox}
{\sffamily\bfseries\small\color{Deep}Two very large paediatric studies land on the same day
and both report that the constituent, not the mass, carries the effect --- and both
identify \emph{organic matter} as the dominant contributor.}\\[1.4mm]
\footnotesize
Wu et al.\ (\textit{Ecotoxicol Environ Saf}) run quantile g-computation on
\textbf{339{,}435 mother--child pairs} in Liaoning and find prenatal PM$_{2.5}$ total mass
at \textbf{OR 1.050 (95\% CI 1.036--1.064) per IQR} for overweight/obesity by age two,
while the constituent mixture reaches \textbf{1.063 (1.050--1.076)} with
\textbf{organic matter contributing 54.1\%} of the joint effect, ammonium 34.4\% and
nitrate 11.5\%. Wang et al.\ (same journal, \textbf{449{,}894 children under five} in
China, India and 26 African countries) report cumulative PM$_{2.5}$ at
\textbf{OR 1.019 (1.009--1.028) in China, 1.009 (1.008--1.009) in India and
1.001 (1.000--1.002) in Africa} for anaemia, with organic carbon and dust the leading
components and population attributable fractions of \textbf{28.5\%, 10.6\% and 0.69\%}
respectively. The two papers disagree on almost everything else --- design, outcome,
exposure product (Chinese reanalysis vs.\ MERRA-2), continent --- and converge on the same
methodological demand: \emph{a mass-only instrument cannot measure the quantity that
carries the association}. For a distributed low-cost network this is the recurring problem
in its sharpest form, because optical scattering is precisely the technique least able to
separate organic matter from ammonium nitrate. Note also that the PAF ordering is inverted
relative to the exposure ordering --- Africa has the highest anaemia prevalence (64.5\%)
and the lowest attributable fraction --- which is a statement about model-assigned exposure
contrast, not about African air.
\end{keybox}

\vspace{3mm}

% -------------------------------------------------------------- exec brief
\section*{\sffamily\bfseries\color{Deep}What changed today}
\vspace{-1.5mm}{\color{Amber}\rule{\linewidth}{1.1pt}}\vspace{2.5mm}

\footnotesize
\begin{itemize}
\item \textbf{Effect modification, not main effect, is where the COVID--PM$_{2.5}$ signal
now sits.} Kannoth et al.\ (\textit{JESEE}) find short-term PM$_{2.5}$ alone essentially
null for COVID-19 pneumonia (\textbf{RR 1.03, 95\% CI 0.95--1.12}) and prior myocardial
infarction alone at \textbf{1.34 (1.14--1.58)}, but the joint contrast at
\textbf{1.67 (1.41--1.99)}, with upper-quartile PM$_{2.5}$ raising pneumonia risk
\emph{only} in those with cardiovascular disease (ATE 0.05, 0.02--0.08). A study powered on
the main effect would have concluded nothing.

\item \textbf{Non-exhaust PM gets a validated inorganic workflow.} Diana et al.\
(\textit{Talanta}) put 32 WLTP brake-dynamometer PM$_{10}$ samples through a dual digestion
(HNO$_3$/H$_2$O$_2$ and HNO$_3$/HF/H$_3$BO$_3$) validated against NIST SRM 1648a,
quantifying \textbf{27 elements at $>$75\% recovery}. Low-steel pads on grey cast iron give
the highest emission factors and Fe-dominated PM$_{10}$; non-asbestos organic pads give the
lowest and are marked by Ba, Ti, K and Zr. That is a usable source-marker set for
apportionment as tailpipe emissions keep falling and brake wear does not.

\item \textbf{Two independent attacks on the same instrumentation problem: extracting
structure from sparse regulatory data.} Peri\v{s}i\'c et al.\ (\textit{Sci Total Environ})
cluster SHAP profiles from a single station into ten recurrent
meteorology--pollutant settings ($R^2 = 0.913$), showing that identical PM$_{10}$
concentrations correspond to different physical configurations. Band \& Mohammadi (arXiv)
pretrain on 10 US stations and adapt to 77 in Taiwan, reaching \textbf{$R^2 = 0.8739$,
MSE 21.66, MAE 3.14} with an adaptive source encoder. Both are aimed at deployments that
cannot afford speciation or dense reference coverage.

\item \textbf{The day's health records are all exposure-modelled, none exposure-measured.}
Every one of the four human-endpoint records assigns exposure from a reanalysis or CMAQ
product at ZIP-code, district or 1\,km resolution. No record this issue measured particles
and outcomes on the same subjects --- the third consecutive issue with that property.

\item \textbf{Weak record worth flagging rather than hiding.} Erickson-DiRenzo et al.\
(\textit{J Voice}) report paired laryngeal biopsies from \textbf{one} participant during a
117\,$\mu$g/m$^3$ wildfire episode and at $<$10\,$\mu$g/m$^3$, with a $\approx$2-fold
MMP-1 increase. $n=1$, no control for time, phonatory load or infection; carried for the
paired-biopsy design, not the result.
\end{itemize}

\clearpage

% ==================================================================== ANALYTICS
\section*{\sffamily\bfseries\color{Deep}Corpus analytics}
\vspace{-1.5mm}{\color{Amber}\rule{\linewidth}{1.1pt}}\vspace{3mm}

\noindent\includegraphics[width=\linewidth]{\FIGDIR/f1_subtopics.png}
\figcap{Subtopic assignment of the 10 carried records, single-label by dominant
contribution. Ten records is the smallest issue of the watch to date, tying 5 August; the
window is a single Sunday, and PubMed's \texttt{[EDAT]} pipeline deposits almost nothing at
the weekend --- the connector returned one record for the whole day. Six of ten clusters
are occupied. \textit{Sensing, forecasting \& instrumentation} leads at 3;
\textit{Neuro / mental health}, \textit{Cardiovascular \& metabolic},
\textit{Mechanistic toxicology} and \textit{Occupational \& indoor} are all empty. Neuro is
now empty on 11 of 22 issues and remains the watch's least reliable cluster.}

\vspace{2mm}
\noindent\includegraphics[width=\linewidth]{\FIGDIR/f2_design_endpoint.png}
\figcap{Left --- study architecture. \textit{Modelling / inventory} takes 4 of 10, the
largest single block: the XAI station study, the transfer-learning forecaster, the
nucleation cluster calculation and the Riyadh satellite preprint. Observational designs
split 2 cohort, 1 cross-sectional, 1 ecological; one measurement campaign (the paired
laryngeal biopsies) and one bench experiment (the brake dynamometer). No
\textit{Trial / intervention} and no \textit{Review / synthesis} record at all. Right ---
health endpoint. \textbf{No health endpoint} is the largest bar at 6 of 10 (60\%), within
the 41--67\% band the last eight issues have occupied once the
\texttt{None (\ldots)} label family is correctly canonicalised --- see the provenance note.}

\vspace{2mm}
\noindent\includegraphics[width=\linewidth]{\FIGDIR/f3_metric_geography.png}
\figcap{Left --- particle metric. PM$_{2.5}$-only accounts for 5 records; one is
PM$_{10}$-only (the XAI station study), one is PM$_{2.5}$ and PM$_{10}$ jointly, and two
are unspeciated PM or emissions. The single composite-index slice is the Riyadh preprint,
which carries \emph{no surface PM product at all} --- only TROPOMI trace gases and an
Absorbing Aerosol Index; read that slice as an absence. Right --- geography. Europe, China
and North America hold 2 each; East Asia (ex-China), Southeast Asia, MENA and one
multi-region record take the remainder. The multi-region record (Wang et al.) alone spans
28 countries, so the count under-represents the day's actual geographic reach.}

\vspace{2mm}
\noindent\includegraphics[width=\linewidth]{\FIGDIR/f6_lifecourse.png}
\figcap{Life-course windows addressed. Counts are non-exclusive and describe this issue
only. Preconception/in-utero and infancy/early childhood hold 2 each --- the two large
paediatric constituent studies contribute to both, since each assigns a prenatal exposure
window to a postnatal outcome. Puberty/adolescence and 65+ are empty. Working age holds
2: the NYC hospitalisation cohort and the single-participant laryngeal study.}

\vspace{2mm}
\noindent\includegraphics[width=\linewidth]{\FIGDIR/f4_heatmap.png}
\figcap{Subtopic $\times$ study-architecture cross-tabulation. Ten records occupy 8 of the
70 available cells and no cell holds more than 2. With a corpus this small the figure
carries little information beyond the silo already visible in f2: every
\textit{Sensing} and \textit{Exposure} record sits in a \textit{Modelling / inventory},
\textit{Measurement campaign} or \textit{Experimental} column, every human-endpoint record
in an observational one, and the two never meet.}

\vspace{2mm}
\noindent\includegraphics[width=\linewidth]{\FIGDIR/f5_forest.png}
\figcap{The 11 interval estimates carried today, on a common logarithmic axis --- more
estimates than records, because three studies each contribute a family of nested
contrasts. \textbf{The estimates are not mutually independent.} The three Kannoth
contrasts (PM$_{2.5}$ alone, myocardial infarction alone, both jointly) are drawn from the
same 15{,}361-admission cohort and share person-time; the five Wu constituent estimates are
IQR-scaled contrasts on correlated constituents of the same mixture in the same 339{,}435
pairs, so their ordering is informative but their spread is not a between-study spread; the
three Wang estimates are region-stratified from one pooled cross-section. Exposure
contrasts also differ in kind --- Kannoth reports an IQR contrast on a 7-day average, Wu an
IQR contrast on a gestational mean, Wang a per-1\,$\mu$g/m$^3$ slope on a cumulative
window --- so interval width across the three families is not comparable and the plot
should be read within family, not across it.}

\clearpage

% ==================================================================== DIGEST
\section*{\sffamily\bfseries\color{Deep}Paper-level digest}
\vspace{-1.5mm}{\color{Amber}\rule{\linewidth}{1.1pt}}\vspace{1mm}

{\fontsize{7.4}{9}\selectfont\color{Slate}Ordered by cluster. Each entry gives the
methodological core and the finding that matters, not an abstract paraphrase.\par}

\band{Sensing, forecasting \& instrumentation\hfill 3 records}{Amber}

\paperentry{Peri\v{s}i\'c et al.\ 2026}
{Explainable AI resolves ten recurrent environmental settings behind urban PM$_{10}$ variability}
{Sci Total Environ}
{Four years (2020--2023) of routine regulatory data from a \emph{single} station ---
PM$_{10}$, PM$_{2.5}$, NO$_2$, SO$_2$, O$_3$ and meteorology --- fed to ensemble
machine-learning models with metaheuristic hyperparameter optimisation, reaching
\textbf{$R^2 = 0.913$}, then clustered on SHAP-derived predictor-impact profiles rather
than on concentrations. \textbf{Ten recurrent settings emerge.} The strongest positive
contributions are cold-season accumulation regimes: E0 is frequent and predominantly
nocturnal at a mean impact of \textbf{50.9\,$\mu$g\,m$^{-3}$}; E7 is rarer but reaches
\textbf{82.9\,$\mu$g\,m$^{-3}$} and persists day and night. Warm, well-mixed settings run
\textbf{26--29\,$\mu$g\,m$^{-3}$ below} the model baseline. The substantive claim --- that
equal concentrations can encode different pollutant--meteorology configurations --- is well
supported and directly useful where speciation is unaffordable. The limitation is
structural and the authors say so: one station, one climate, no external validation, and
SHAP attributions are model-conditional, so the ``settings'' are properties of a fitted
ensemble before they are properties of the atmosphere. For distributed low-cost networks
the appeal is that the entire input set is what a modest node plus a met sensor already
produces.}
{10.1016/j.scitotenv.2026.182195}

\paperentry{Diana et al.\ 2026}
{Validated multi-technique workflow for inorganic characterisation of brake-wear PM$_{10}$}
{Talanta}
{Thirty-two PM$_{10}$ samples from WLTP-based brake-dynamometer tests on four tribological
couples --- low-steel (LS) and non-asbestos organic (NAO) pads against grey cast-iron (GCI)
or coated discs. A combined acid digestion (HNO$_3$/H$_2$O$_2$ plus
HNO$_3$/HF/H$_3$BO$_3$) was validated against \textbf{NIST SRM 1648a} and quantified by
ICP-OES and ICP-MS, giving \textbf{27 elements with recoveries above 75\% for most
analytes}; FESEM-EDS confirmed the trends at particle scale and PCA/HCA separated the
tribological families cleanly. \textbf{LS-GCI produced the highest PM$_{10}$ emission
factors and NAO-GCI the lowest}; Fe dominates LS-derived material while Ba, Ti, K and Zr
mark NAO. An Mn-based enrichment factor improved discrimination of rotor- versus
pad-derived contributions. The value here is a defensible marker set for non-exhaust source
apportionment, which matters more each year as tailpipe PM falls. Caveats are the usual
dynamometer ones: WLTP braking is not real-world braking duty, only four couples were
tested, and no on-road or roadside validation is attempted.}
{10.1016/j.talanta.2026.130448}

\paperentry{Band \& Mohammadi 2026 (preprint)}
{Shift-aware dual-encoder transfer learning for short-horizon PM$_{2.5}$ forecasting}
{arXiv}
{A source encoder pretrained on hourly PM$_{2.5}$ from 10 US monitoring locations is fused
with a target-specific encoder and adapted to \textbf{77 Taiwanese stations} over two years
under a chronological train/validation/test split --- the correct protocol, and one many
comparable papers do not use. Against four baselines the frozen-source dual encoder gives
\textbf{MSE 21.8960, MAE 3.1597, $R^2 = 0.8725$} (a 7.1\% and 4.1\% MSE reduction over two
transfer-learning variants); allowing the source encoder to adapt improves this to
\textbf{MSE 21.6575, MAE 3.1383, $R^2 = 0.8739$}. Ablation shows the Taiwan-specific branch
matters most, and SHAP attributes predictions mainly to recent PM$_{2.5}$ and meteorology.
The framing --- explicit handling of source/target distribution shift to avoid negative
transfer --- is the right one for anyone standing up a network in a region with no long
record. Two reservations: the absolute gain over the better baselines is a few percent of
MSE and no significance or repeated-seed variance is reported, and both domains are
reference-grade regulatory data, so transfer onto low-cost nodes (a far larger shift) is
untested. Preprint, not peer reviewed.}
{10.48550/arXiv.2608.14456}

\band{Exposure assessment \& modelling\hfill 2 records}{Amber}

\paperentry{Fan et al.\ 2026}
{Multiple amides synergistically enhance sulfuric acid--ammonia nucleation}
{iScience}
{Quantum-chemical cluster thermodynamics plus kinetic simulation of new particle formation
in binary, ternary and quaternary systems containing sulfuric acid (SA), ammonia,
methylacetamide (MAA) and methylpropionamide (MPA). \textbf{(SA)$_1\cdot$(MPA)$_1$ and
(SA)$_1\cdot$(MAA)$_1$ are more stable than (SA)$_1\cdot$(NH$_3$)$_1$ and even than
(SA)$_1\cdot$(DMA)$_1$}, and the two amides together enhance nucleation by
\textbf{an order of magnitude more than either alone}, most strongly at low SA
concentration and low temperature. A displacement hierarchy is identified: MPA replaces
NH$_3$, MAA replaces MPA. Carried because secondary particle number --- not mass --- is
what an optical low-cost sensor is blindest to, and because amide sources (industrial,
agricultural, carbon-capture solvent degradation) are rising and largely unmonitored. It is
a purely computational result with no chamber or ambient confirmation, the amide
concentrations assumed are not tied to measured ambient values, and cluster-level
thermodynamics has a poor track record of quantitative transfer to observed formation
rates.}
{10.1016/j.isci.2026.117013}

\paperentry{Alshehri \& El-Hamid 2026 (preprint)}
{Sentinel-5P and MODIS land-surface coupling of pollutant fields over Riyadh, 2018--2025}
{Research Square}
{Eight years of TROPOMI NO$_2$, SO$_2$, CO, O$_3$ and Absorbing Aerosol Index combined with
MODIS land-surface temperature and NDVI and Sentinel-2 land cover in Google Earth Engine,
analysed by correlation, PCA and multiple linear regression. NO$_2$ and CO concentrate over
urban and industrial areas, O$_3$ peaks downwind, LST falls from 43.10\,$^\circ$C (2021) to
37.07\,$^\circ$C (2025) while NDVI rises after 2022. \textbf{The regression is weak:
$R^2$ ranges from 0.124 (O$_3$) to 0.231 (SO$_2$)}, and PCA explains 55.8\% of variance;
the authors correctly conclude that land-surface descriptors explain only a minor share and
that emissions and transport dominate. Carried as a negative methodological result for an
arid-region setting where ground networks are sparse --- and as a caution, because
\textbf{no surface PM product appears anywhere in the study} despite the framing as air
pollution monitoring; AAI is a column absorption index, not a PM concentration. Preprint,
not peer reviewed, and the LST trend is reported without any test for its significance
against interannual variability.}
{10.21203/rs.3.rs-10538140/v1}

\band{Respiratory \& allergic\hfill 2 records}{Sage}

\paperentry{Kannoth et al.\ 2026}
{Short-term PM$_{2.5}$ and COVID-19 pneumonia in New York City, amplified by pre-existing cardiovascular disease}
{J Expo Sci Environ Epidemiol}
{\textbf{15{,}361 NYC COVID-19 hospitalisations} (March--December 2020) with short-term
PM$_{2.5}$ defined as the 7-day residential ZIP-code average from EPA CMAQ, and 15
ICD-coded circulatory conditions as the modifier. Modified Poisson regression for ARDS,
pneumonia, ventilation and dialysis; Cox models for fatality and length of stay; RERI for
additive interaction; TMLE for average treatment effects. \textbf{PM$_{2.5}$ alone was null
(RR 1.03, 95\% CI 0.95--1.12) and prior myocardial infarction alone 1.34 (1.14--1.58), but
jointly 1.67 (1.41--1.99)}; upper-quartile PM$_{2.5}$ raised pneumonia risk only among
those with cardiovascular conditions (\textbf{ATE 0.05, 0.02--0.08}). The interaction is
the finding and it is the sort a main-effect analysis destroys. Weaknesses are real:
ZIP-code CMAQ is a coarse surrogate for personal exposure and its error is spatially
structured with deprivation, the effect appears only from July onward with no mechanism
offered for the seasonal split, and the cohort is hospitalised patients only, so selection
on admission is unaddressed. For network design the implication is concrete --- the
susceptible stratum is identifiable in advance, which is an argument for siting density by
cardiovascular prevalence rather than by population alone.}
{10.1038/s41370-026-00957-5}

\paperentry{Erickson-DiRenzo et al.\ 2026}
{Wildfire PM$_{2.5}$ and the human larynx: paired in-event and clean-air assessment in one adult}
{J Voice}
{A single adult assessed twice --- during a wildfire episode with PM$_{2.5}$ reaching
\textbf{117\,$\mu$g\,m$^{-3}$} and again at \textbf{$<$10\,$\mu$g\,m$^{-3}$} --- with
acoustic and aerodynamic voice measures, laryngeal endoscopy with stroboscopy, and
\emph{ventricular fold biopsies at both time points} analysed by qPCR for extracellular
matrix and inflammatory markers. Speaking fundamental frequency fell relative to baseline,
endoscopy showed visibly increased laryngeal mucus, and \textbf{MMP-1 expression roughly
doubled}; other ECM and inflammatory genes moved inconsistently. \textbf{$n = 1$.} There is
no control participant, no adjustment for phonatory load, hydration, intercurrent
infection or the ordinary diurnal variation in fundamental frequency, and a two-point
comparison cannot separate exposure from time. The reason to carry it is the design, not
the result: paired human airway tissue across a natural exposure contrast is rare, and it
establishes feasibility for a properly powered version. Treat the MMP-1 figure as a
hypothesis.}
{10.1016/j.jvoice.2026.07.048}

\band{Reproductive \& developmental\hfill 1 record}{Sky}

\paperentry{Wu et al.\ 2026}
{Prenatal PM$_{2.5}$ constituents and early childhood obesity: organic matter dominates the mixture}
{Ecotoxicol Environ Saf}
{\textbf{339{,}435 mother--child pairs} in Liaoning Province with prenatal PM$_{2.5}$ and
five constituents (organic matter, black carbon, ammonium, nitrate, sulfate) assigned from
high-resolution atmospheric reanalysis, and overweight/obesity ascertained across the first
two years. Single-pollutant models give \textbf{PM$_{2.5}$ mass OR 1.050 (95\% CI
1.036--1.064) per IQR}, with \textbf{OM 1.058 (1.043--1.072)}, \textbf{NH$_4^+$ 1.048
(1.035--1.062)}, \textbf{NO$_3^-$ 1.048 (1.033--1.064)} and \textbf{BC 1.033
(1.020--1.047)}; sulfate is null. Quantile g-computation puts the mixture at \textbf{1.063
(1.050--1.076)} per quartile increase with \textbf{OM contributing 54.1\%}, NH$_4^+$ 34.4\%
and NO$_3^-$ 11.5\%. The sample size makes the intervals tight, but tight intervals on
reanalysis-assigned exposure are a precision claim about the model, not the exposure:
constituent fields are far more uncertain than total mass and their errors are correlated
with each other, which is exactly the condition under which QGC weights become unstable.
Residential address at delivery stands in for the whole gestation, and no gestational
windows are resolved. The direction is credible; the 54\% attribution to OM is not
something the exposure product can really support.}
{10.1016/j.ecoenv.2026.120670}

\band{Burden, policy \& mitigation\hfill 1 record}{Coral}

\paperentry{Alamsyah et al.\ 2026}
{Industrial intensity and traffic congestion propagate air-pollution health vulnerability across district borders in Java}
{Environ Sci Pollut Res Int}
{A health vulnerability index (HVI) for Java built by PCA over satellite pollutant fields,
official statistics and OpenStreetMap congestion data, then regressed on industrial
activity intensity and congestion in a spatial Durbin model to separate direct from
spillover effects. \textbf{HVI ranges 29.2--62.3 on a 0--100 scale}, highest in West Java,
Jakarta and Banten and lowest in Central Java, Yogyakarta and East Java; the exposure
dimension varies most (13.18) against adaptive capacity (9.40) and sensitivity (5.17), and
both industrial intensity and congestion act through \textbf{direct and indirect
(spillover) channels}. The spillover result is the useful one --- it says district-level
mitigation will leak, which is a real argument for regional rather than municipal air
management. But the index is a construct: PCA weights are data-driven rather than
health-anchored, no health outcome is used to validate the HVI against observed morbidity,
and OpenStreetMap congestion is an availability-biased proxy that is denser precisely in
the wealthier districts. Uncertainty and sensitivity analysis is reported, which is more
than most index papers do.}
{10.1007/s11356-026-38148-2}

\band{Other clinical endpoints\hfill 1 record}{Slate}

\paperentry{Wang et al.\ 2026}
{Early-life PM$_{2.5}$ mass and components and childhood anaemia across China, India and Africa}
{Ecotoxicol Environ Saf}
{Cross-sectional analysis of \textbf{449{,}894 children under five} from China, India and 26
African countries, with prenatal, postnatal and cumulative PM$_{2.5}$ and five components
(black carbon, organic carbon, sulfate, dust, sea salt) from NASA MERRA-2 reanalysis.
Anaemia prevalence is \textbf{26.9\% in China, 62.6\% in India, 64.5\% in Africa}. Per
1\,$\mu$g\,m$^{-3}$ of cumulative PM$_{2.5}$: \textbf{OR 1.019 (95\% CI 1.009--1.028) in
China, 1.009 (1.008--1.009) in India, 1.001 (1.000--1.002) in Africa}, with organic carbon
and dust positive in all three. Population attributable fractions are \textbf{28.46\%,
10.64\% and 0.69\%}. The regional ordering of effect size and PAF runs \emph{opposite} to
the ordering of anaemia burden, and the paper does not resolve this: the most economical
explanation is that MERRA-2 at 0.5$^\circ\times$0.625$^\circ$ has little within-country
exposure contrast in Africa and correspondingly attenuated coefficients, not that African
children are less susceptible. Cross-sectional design cannot order exposure and outcome,
and anaemia in these settings is dominated by iron deficiency, malaria and haemoglobinopathy
--- confounders that co-vary strongly with dust exposure. Read the component ranking, not
the PAFs.}
{10.1016/j.ecoenv.2026.120658}

\clearpage

% ==================================================================== TABLE
\section*{\sffamily\bfseries\color{Deep}Corpus register}
\vspace{-1.5mm}{\color{Amber}\rule{\linewidth}{1.1pt}}\vspace{2.5mm}

{\fontsize{6.7}{8.4}\selectfont
\setlength{\tabcolsep}{2.4pt}
\begin{longtable}{@{}>{\RaggedRight}p{25mm}>{\RaggedRight}p{41mm}>{\RaggedRight}p{27mm}%
>{\RaggedRight}p{23mm}>{\RaggedRight}p{18mm}>{\RaggedRight}p{15mm}@{}}
\toprule
\textsf{\textbf{First author}} & \textsf{\textbf{Focus}} & \textsf{\textbf{Journal}} &
\textsf{\textbf{Design}} & \textsf{\textbf{Metric}} & \textsf{\textbf{Region}}\\
\midrule
\endfirsthead
\toprule
\textsf{\textbf{First author}} & \textsf{\textbf{Focus}} & \textsf{\textbf{Journal}} &
\textsf{\textbf{Design}} & \textsf{\textbf{Metric}} & \textsf{\textbf{Region}}\\
\midrule
\endhead
\multicolumn{6}{@{}l@{}}{\cellcolor{Mist}\textsf{\textbf{\textcolor{Deep}{Neuro / mental health}}}}\\[0.6mm]
\href{https://doi.org/10.1001/jamanetworkopen.2026.25015}{Dales} & Wildfire-sourced PM$_{2.5}$ and ED visits for migraine and other primary headache & \textit{JAMA Netw Open} & Case-crossover & PM$_{2.5}$ (wildfire) & Canada\\
\href{https://doi.org/10.1093/ageing/afag215}{Glover} & Umbrella meta-review of urban environmental exposures and cognitive health & \textit{Age Ageing} & Umbrella review / meta-analysis & PM$_{2.5}$ & Global\\
\href{https://doi.org/10.1016/j.nsa.2026.107019}{Goossens} & In utero air pollution exposure and pre-/postnatal brain development & \textit{Neurosci Appl} & Narrative review & PM$_{2.5}$ / PM$_{10}$ / PM$_{\mathrm{coarse}}$ & Global\\
\href{https://doi.org/10.1016/j.isci.2026.116617}{Ma} & PM$_{2.5}$-induced neurotoxicity and depression: exposure risks and mechanisms & \textit{iScience} & Narrative review & PM$_{2.5}$ + components & Global\\
\href{https://doi.org/10.3760/cma.j.cn112150-20250804-00752}{Pan} & Maternal PM exposure and offspring neurodevelopment: effects and mechanisms & \textit{Zhonghua Yu Fang Yi Xue Za Zhi} & Narrative review & PM (bulk) & China\\
\addlinespace[0.9mm]
\multicolumn{6}{@{}l@{}}{\cellcolor{Mist}\textsf{\textbf{\textcolor{Teal}{Cardiovascular \& metabolic}}}}\\[0.6mm]
\href{https://doi.org/10.3760/cma.j.cn112338-20251120-00843}{Huang} & PM$_{2.5}$ and dyslipidaemia in rural Chinese children; dietary effect modification & \textit{Zhonghua Liu Xing Bing Xue Za Zhi} & Cross-sectional (multistage) & PM$_{2.5}$ & China\\
\href{https://doi.org/10.32394/pe/218966}{Kazmierska} & Air pollution and cardiovascular disease: review of current data & \textit{Przegl Epidemiol} & Narrative review & PM$_{2.5}$ / PM$_{10}$ & Global\\
\href{https://doi.org/10.1007/s00125-026-06814-2}{Xi} & Air pollution, built environment and temperature vs. prediabetes phenotypes & \textit{Diabetologia} & Prospective cohort & PM$_{2.5}$ & Germany\\
\href{https://doi.org/10.1161/JAHA.125.046917}{Yang} & Short-term PM$_{2.5}$ and hospitalisation/ED use in atrial fibrillation & \textit{J Am Heart Assoc} & Case-crossover & PM$_{2.5}$ (wildfire-influenced) & USA\\
\addlinespace[0.9mm]
\multicolumn{6}{@{}l@{}}{\cellcolor{Mist}\textsf{\textbf{\textcolor{Sky}{Reproductive \& developmental}}}}\\[0.6mm]
\href{https://doi.org/10.1016/j.jes.2025.10.043}{Chen} & Prenatal air pollution, fetal growth trajectory phenotypes and child neurodevelopment & \textit{J Environ Sci (China)} & Prospective cohort & PM$_{2.5}$ / PM$_{10}$ & China\\
\href{https://doi.org/10.1016/j.ecoenv.2026.120566}{Huang} & Joint prenatal bisphenol + air-pollutant mixture exposure and preschool IQ & \textit{Ecotoxicol Environ Saf} & Prospective cohort (mixtures) & PM$_{2.5}$ / PM$_{10}$ & China\\
\href{https://doi.org/10.1007/s40572-026-00552-8}{Peterson} & Ambient air pollution, PAHs and pubertal development: critical review & \textit{Curr Environ Health Rep} & Critical review & PM$_{2.5}$ components & Global\\
\href{https://doi.org/10.1016/j.envres.2026.125313}{Xie} & Long-term PM$_{2.5}$ and timing of pubertal development in boys and girls & \textit{Environ Res} & Prospective cohort & PM$_{2.5}$ & China\\
\addlinespace[0.9mm]
\multicolumn{6}{@{}l@{}}{\cellcolor{Mist}\textsf{\textbf{\textcolor{Sage}{Respiratory \& allergic}}}}\\[0.6mm]
\href{https://doi.org/10.1007/s00484-026-03289-z}{Bayar Muluk} & Climate change and otolaryngology: airway disease impacts and emission reduction & \textit{Int J Biometeorol} & Narrative review & PM$_{2.5}$ & Global\\
\href{https://doi.org/10.4103/mgr.MEDGASRES-D-25-00164}{Liao} & Short-term air pollution, respiratory symptoms and lung function in asthma & \textit{Med Gas Res} & Panel study & PM$_{2.5}$ / PM$_{10}$ & China\\
\href{https://doi.org/10.3390/epidemiologia7040097}{Sanchez de Prada} & Ambient PM and asthma-coded admissions in children 0-3 y, Spain & \textit{Epidemiologia} & Ecological panel & PM$_{2.5}$ / PM$_{10}$ & Spain\\
\href{https://doi.org/10.1016/j.jhazmat.2026.143076}{Yang} & Allergenic vs. toxic risk of airborne seed hairs compared with ambient PM$_{2.5}$ & \textit{J Hazard Mater} & Field measurement & PM$_{2.5}$ + bioaerosol & China\\
\addlinespace[0.9mm]
\multicolumn{6}{@{}l@{}}{\cellcolor{Mist}\textsf{\textbf{\textcolor{Violet}{Mechanistic toxicology}}}}\\[0.6mm]
\href{https://doi.org/10.1016/j.jes.2025.12.074}{Ding} & PM$_{2.5}$ promotes vascular calcification via miR-27b-5p/HES1 signalling & \textit{J Environ Sci (China)} & In vitro + organoid + murine & PM$_{2.5}$ & China\\
\href{https://doi.org/10.1007/s00360-026-01698-5}{Maraschi} & Osmoregulatory disruption and metal bioaccumulation from settleable atmospheric PM & \textit{J Comp Physiol B} & Controlled exposure (ecotox) & Settleable APM & Brazil\\
\href{https://doi.org/10.1016/j.jes.2025.11.018}{Wu} & Macrophage NLRP3 inflammasome in real-ambient PM-induced hepatic glucose dysregulation & \textit{J Environ Sci (China)} & Real-ambient murine + co-culture & Real-ambient PM & China\\
\addlinespace[0.9mm]
\multicolumn{6}{@{}l@{}}{\cellcolor{Mist}\textsf{\textbf{\textcolor{Clay}{Occupational \& indoor}}}}\\[0.6mm]
\href{https://doi.org/10.3389/fmedt.2026.1796828}{Chen} & Nested suction instrument to cut intraoperative surgical-smoke PM$_{2.5}$ & \textit{Front Med Technol} & Bench / device experiment & PM$_{2.5}$ (surgical smoke) & China\\
\href{https://doi.org/10.1093/annweh/wxag055}{Cole} & Determinants of respirable crystalline silica exposure in Australian tunnelling & \textit{Ann Work Expo Health} & Observational exposure study & Respirable dust / RCS & Australia\\
\href{https://doi.org/10.1016/j.jhazmat.2026.143081}{Ding} & Stage-resolved hair Zn-isotope dynamics as a firefighting combustion-exposure tracer & \textit{J Hazard Mater} & Repeated-measures biomarker & Combustion-derived PM & China\\
\href{https://doi.org/10.3390/jox16040119}{Ramsden} & Xenobiotic hazards in aircraft cabin air & \textit{J Xenobiot} & Review + risk economics & Ultrafine particles & Global\\
\addlinespace[0.9mm]
\multicolumn{6}{@{}l@{}}{\cellcolor{Mist}\textsf{\textbf{\textcolor{Amber}{Sensing, forecasting \& instrumentation}}}}\\[0.6mm]
\href{https://doi.org/10.1016/j.jenvman.2026.130515}{Shi} & Adaptive spatiotemporal graph transformer for multi-site multi-step PM$_{2.5}$ forecasting & \textit{J Environ Manage} & Deep-learning model & PM$_{2.5}$ & China / India\\
\href{https://doi.org/10.3390/bios16070380}{Tsuchiya} & Seasonal and indoor-outdoor characteristics of fluorescent aerosol particles in offices & \textit{Biosensors} & Real-time sensor campaign & Bioaerosol / coarse & Japan\\
\addlinespace[0.9mm]
\multicolumn{6}{@{}l@{}}{\cellcolor{Mist}\textsf{\textbf{\textcolor{Amber}{Exposure assessment \& modelling}}}}\\[0.6mm]
\href{https://doi.org/10.1029/2025GH001621}{Solorzano-Araque} & Multiscale exposure modelling: LOTOS-EUROS + mobility surveys + particle toxicity & \textit{GeoHealth} & CTM + mobility integration & PM$_{2.5}$ / PM$_{10}$ & Colombia\\
\addlinespace[0.9mm]
\multicolumn{6}{@{}l@{}}{\cellcolor{Mist}\textsf{\textbf{\textcolor{Coral}{Burden, policy \& mitigation}}}}\\[0.6mm]
\href{https://doi.org/10.1002/hsr2.72886}{Karamzad} & Ischaemic heart disease burden in MENA 1990-2021; PM attributable fraction 28.2\% & \textit{Health Sci Rep} & GBD secondary analysis & Ambient PM & MENA\\
\href{https://doi.org/10.1371/journal.pone.0347989}{Rees} & HEPA filtration to reduce respiratory infection in care homes & \textit{PLoS One} & RCT process evaluation & Filtered PM & UK\\
\href{https://doi.org/10.1016/j.jes.2026.01.009}{Shan} & Policy-driven emission reductions and post-pandemic rebound, high-resolution inventories & \textit{J Environ Sci (China)} & Emission inventory & PM emissions & China\\
\addlinespace[0.9mm]
\multicolumn{6}{@{}l@{}}{\cellcolor{Mist}\textsf{\textbf{\textcolor{Slate}{Other clinical endpoints}}}}\\[0.6mm]
\href{https://doi.org/10.3390/dj14070440}{Lupita} & PM$_{2.5}$, preventive-care capacity and caries/edentulism burden across the EU & \textit{Dent J} & Ecological panel & PM$_{2.5}$ & EU-27\\
\href{https://doi.org/10.1016/j.jdin.2026.06.008}{Park} & Long-term PM exposure and risk of androgenetic alopecia & \textit{JAAD Int} & Nationwide cohort & PM$_{2.5}$ / PM$_{10}$ & South Korea\\
\href{https://doi.org/10.3171/2026.3.PEDS2615}{Sun} & Environmental exposure and area deprivation vs. paediatric hydrocephalus outcomes & \textit{J Neurosurg Pediatr} & Registry cohort + geospatial & PM$_{2.5}$ (satellite-derived) & USA\\
\addlinespace[0.9mm]
\bottomrule
\end{longtable}}


\vspace{4mm}
\begin{keybox}
{\sffamily\bfseries\footnotesize\color{Deep}Provenance and method}\\[1.2mm]
{\fontsize{7.2}{8.8}\selectfont
\textbf{Window.} Entry window \textbf{2026-08-16 to 2026-08-16}, contiguous with the
15 August issue; no gap. Every record was checked against \texttt{state/seen.json}
(pmid $\rightarrow$ lowercased doi $\rightarrow$ title signature) before summarising;
\textbf{11 Europe PMC hits were dropped as already-seen duplicates}.\\[0.8mm]
\textbf{Harvester defect from 15 August: repaired.} The Europe PMC leg in
\texttt{harvest.py} was windowing on \texttt{FIRST\_PDATE}, which the source-status notes
had already recorded as non-functional for recent windows --- it returned 0 for both the
15 and 16 August windows. Switched to \texttt{CREATION\_DATE} and given cursor-mark
pagination to 600 records; the committed leg now returns \textbf{20} for this window
unaided, and every record in this issue except one came through it.\\[0.8mm]
\textbf{Per-leg yield.} A single Sunday window is close to the floor of what this watch can
see. PubMed connector, two separate queries on the health and instrumentation axes:
\textbf{1 hit total}, the same record on both axes, already present in the Europe PMC set.
Local harvester \texttt{pubmed\_health} 0 and \texttt{pubmed\_sensing} 0 --- PubMed's
\texttt{[EDAT]} pipeline deposits almost nothing at the weekend, which is a property of the
source and not a failure. Europe PMC primary query 20 hits, 9 new after deduplication,
\textbf{6 carried}; a second broader Europe PMC sweep on air-quality and aerosol terms
returned 36, 15 further new, \textbf{3 carried}. arXiv 145\,kB with one in-scope entry
dated 14 August that the 15 August sweep did not surface, \textbf{1 carried}. Crossref by
ISSN returned \textbf{0} across 18 journals. \textbf{OpenAlex HTTP 429, 14th consecutive}
--- the leg is dead. \textbf{Consensus: 8th consecutive zero-carry.} All three returned
records were in scope and useful (AutoML indoor calibration, multi-city bi-pollutant LCS
calibration, global PurpleAir transferability) but their Crossref \texttt{created} dates
are 2026-01-23, 2026-01-20 and 2026-06-15. It ranks by relevance, not recency, and should
be treated as a weekly/monthly gap-filler only. ClinicalTrials.gov: no PM or air-pollution
registration created or updated inside the window; \texttt{trials.json}
unchanged.\\[0.8mm]
\textbf{Second defect repaired: endpoint canonicalisation.} \texttt{ENDPOINT\_CANON} mapped
\texttt{None (monitoring)} but not the rest of the \texttt{None (\ldots)} family ---
\texttt{None (exposure)}, \texttt{None (instrument)}, \texttt{None (emissions)},
\texttt{None (algorithm)}, \texttt{None (modelling)}, \texttt{None (aerosol chemistry)} and
singular variants, \textbf{35 records across the store}. All mean the same thing, so f2b
understated its own largest bar; the 9, 10 and 11 August issues rendered a literal
\textbf{0} for a category that in fact held 8, 9 and 15 records. \texttt{canon\_ep} now
falls back on a prefix rule, and the corrected series for 7--16 August reads
54, 64, 67, 56, 58, 46, 41, 43, 43, 60\,\%. \textbf{The three affected daily PDFs are stale
on that panel} and are left as shipped; the store and all future rollups are
correct.\\[0.8mm]
\textbf{Rejections (15).} Six keyword collisions where \texttt{aerosol} or
\texttt{exposure} matched an unrelated field (freeze-spray insecticide delivery,
\textit{Francisella} aerosol-challenge vaccinology, inhaled siRNA therapeutics,
PPE-adherence auditing, lanthanum sediment geochemistry, aqueous biochar sorption); four
gas-phase-only records with no particulate measurement (formaldehyde catalysis,
formaldehyde occupational olfaction, MgAl-LDH NO$_2$ sorption, ozone melanogenesis); a
meteorology-only allergic-rhinitis DLNM carrying pollution as an adjustment covariate; a
therapeutics review naming PM only as a background trigger; two correction/corrigendum
notices; and one journal news item reporting a study already carried on 27 July. All logged
with reasons to \texttt{state/rejected.jsonl}.\\[0.8mm]
\textbf{Two records are preprints} (Band \& Mohammadi, arXiv; Alshehri \& El-Hamid,
Research Square) and are labelled as such in the register and the digest. No record this
issue is metadata-only.\\[0.8mm]
\textbf{Labels.} Subtopic, design, particle-metric and geography labels were drawn only
from vocabularies already mapped in \texttt{plots.py} and \texttt{mktable.py}; no new leaf
label was written this issue and the unmapped-label diagnostics printed nothing.
Assignment is single-label by dominant contribution and is a judgement, not a controlled
vocabulary. \texttt{check\_dois.py} returned \textbf{0 fail, 4 warn} across all 10 records;
all four warnings are title-overlap heuristics firing on editorial focus lines (three) or
on arXiv's absence from Crossref (one), and each was confirmed by eye against the
registered record.\\[0.8mm]
Content derived from PubMed, Europe PMC, Crossref and arXiv metadata; abstracts remain the
copyright of their respective publishers. All DOI links resolve to the publisher of record.\par}
\end{keybox}

\end{document}
